\section{Discussion}

\subsection{Mechanism Validation}

Our experiments validate the complete causal chain underlying error-type gating:
\begin{enumerate}
\item Fine-grained penalties concentrate gradients at traceback locations (H-M1: 16.11$\times$)
\item Traceback accuracy varies dramatically by error type (H-M2: 100\% vs 20\%)
\item Unreliable tracebacks cause gradient noise at ground-truth locations (H-M3: $p<10^{-13}$)
\item Gating preserves signal where accurate, removes noise where not (H-M4: +4.61\%)
\end{enumerate}

The end-to-end efficiency improvement (H-E1) demonstrates practical benefit, though we note this is proof-of-concept scale requiring full validation.

\subsection{Honest Limitations}

\textbf{Synthetic ground truth:} H-M2 and H-M3 use synthetic code templates with predetermined bug locations. While this ensures deterministic validation, real APPS samples may show different accuracy patterns. We view this as a conservative choice---demonstrating the mechanism on controlled samples before applying to noisy real data.

\textbf{Marginal significance for H-M4:} The SNR improvement ($p=0.112$) does not reach conventional significance. Our sample size (200 total) was reduced for PoC scope. The consistent direction across 1000 bootstrap iterations suggests the effect is real but requires $N\approx500$+ for adequate power.

\textbf{Model scale:} All experiments use CodeT5-small (60M parameters). Gradient dynamics may differ at the 770M (CodeT5-large) or 7B+ scale typical of modern code LLMs. However, the causal mechanism---traceback parsing, token-level penalties, gradient concentration---is architecture-agnostic.

\textbf{Untested predictions:} P2 (error distribution shift during training) and P3 (increasing advantage over training) were not tested in our PoC. These require multi-epoch training runs with distribution tracking at checkpoints.

\subsection{Broader Impact}

Our work introduces credit assignment reliability as a lens for understanding feedback granularity. This framing extends beyond error-type gating:
\begin{itemize}
\item \textbf{Test case reliability:} Individual test cases may have different diagnostic value; weighting feedback by test reliability follows similar logic
\item \textbf{Multi-file localization:} Repository-level code generation faces even more complex localization challenges
\item \textbf{Other modalities:} Image/video generation with localized feedback faces analogous reliability questions
\end{itemize}

The key insight---that finer granularity helps only when localization is accurate---applies whenever RL uses spatially localized feedback signals.

\subsection{Future Directions}

\textbf{Continuous weighting:} Binary gating is simple but coarse. Per-exception-type accuracy estimates could enable proportional penalty weighting.

\textbf{Larger scale validation:} Replicating on CodeT5-large and decoder-only models (CodeLlama, DeepSeek-Coder) would establish generality.

\textbf{Real sample annotation:} Developing reliable ground-truth annotation for real APPS samples would strengthen H-M2/H-M3 claims.

\textbf{Distribution dynamics:} Tracking U\_ignore fraction during full training runs would test P2/P3 predictions and potentially reveal time-varying optimal gating strategies.
